% (C) 2026 Ralf Stephan, in collaboration with Claude Code.
% Released under CC0 1.0 Universal (public-domain dedication).
% See https://creativecommons.org/publicdomain/zero/1.0/

\documentclass[11pt]{article}
\usepackage[T1]{fontenc}
\usepackage[utf8]{inputenc}
\usepackage{lmodern}
\usepackage[margin=2.5cm]{geometry}
\usepackage{amsmath,amssymb}
\usepackage{braket}
\usepackage[hidelinks]{hyperref}

% Lean names and paths: verbatim, breakable at . _ /
\DeclareUrlCommand\lean{\urlstyle{tt}}
\def\UrlBreaks{\do\.\do\_\do\/}
% Marks an open problem that blocks checking a main claim of the brief
\newcommand{\blocker}[1]{\textsc{[blocks verification: #1]}~}
% Extra stretch for paragraphs with long, unbreakable Lean names
\emergencystretch=6em

\title{Open problems:\\ Physical vector spaces for quantum mechanics}
\author{}
\date{}

\begin{document}
\maketitle

\noindent
Open problems raised by the technical brief in this directory, and by its Lean formalization
(see \lean{formalization-log.tex}).
Items marked \textsc{[blocks verification]} stand in the way of checking a main
claim of the brief as the brief states it: a \emph{formalization gap} can be closed
in Lean alone, while a \emph{brief issue} needs the brief to be corrected or made
precise first.

\section*{From the brief}
\begin{enumerate}
\item \textbf{Completeness of the Lie algebra.} \emph{Research.} Definition~1
  admits pathological cases: a Hilbert space with the defining observables of a qubit (``too many
  vectors'') and a Schwartz space with position as the only defining observable (``too few
  observables''). ``Addressing this gap is an open problem''; the brief expects a ``final
  definition'' once it is solved.
  Partial answer (\lean{further-results.tex}, Sec.~2): the second case is not a quantum vector
  space, since with position alone $\mathcal{S}(\mathbb{R})$ is not complete
  (\lean{FurtherResults.not_isQuantumVectorSpace_posOnly}). The first case is one
  (\lean{FurtherResults.isQuantumVectorSpace_pair}, for $\mathbb{C}^2 \otimes \mathbb{C}^2$),
  and irreducibility excludes it: every continuous operator that commutes with all defining
  observables is a multiple of the identity (\lean{FurtherResults.not_isIrreducible_pair}).
  This is a form of Roberts's condition (J.\ Math.\ Phys.\ 7 (1966) 1097, p.~1099), and
  Examples~3 and~4 satisfy it (\lean{FurtherResults.isIrreducible_schwartz},
  \lean{FurtherResults.isIrreducible_qubit}). Open: whether irreducibility is the condition
  the brief is looking for, and how it relates to taking the smooth vectors of an irreducible
  unitary group representation as $V$.
  (Brief, Sec.~5, ``Completeness of the Lie algebra''; Sec.~6.)
\item \textbf{Lie algebras determine physical systems.} \emph{Research.} ``Two systems
  are physically equivalent if and only if their Lie algebras of observables are equivalent'' is
  asserted, and is what makes QSTO:V consistent with [B]STDIST. (Brief, Theorem~5, proof; Log,
  ``Differences'', last item.)
  \emph{Related} (\lean{further-results.tex}, Sec.~5). Without the observables, the space alone
  does not tell the systems apart: all $\mathcal S(\mathbb R^n)$ are isomorphic. The
  Gelfand--Shilov spaces $S^{1/2}_{1/2}(\mathbb R^n)$ (exponentially decreasing Hermite
  coefficients) do: they are pairwise non-isomorphic as locally convex spaces (proved there from
  the literature, not formalized). They do not fit Definition~1, however. Their own topology is not
  second countable, and with $X$, $P$, $I$ the space $S^{1/2}_{1/2}(\mathbb R)$ is not complete
  (\lean{FurtherResults.not_isQuantumVectorSpace_gelfandShilov}).
\item \textbf{Ensembles on a quantum vector space.} \emph{Research.} QSTO:V uses
  density operators $\rho$ and $S(\rho) = -\operatorname{tr}(\rho\log\rho)$, and Sec.~5 calls the
  set of ensembles ``the most natural space to consider'', but no ensemble space over $V$ is
  defined. In particular, whether the ensembles induce on the pure states the quotient topology
  of $V$, under which the pure states are second countable exactly when $V$ is, is open.
  \emph{Partial answer} (\lean{further-results.tex}, Sec.~4). For Example~3 the positive
  trace-one Schwartz operators of Keyl, Kiukas and Werner are a natural ensemble space: they are
  the density operators with all moments $\operatorname{tr}\rho N^j$ finite
  (\lean{FurtherResults.IsSchwartzDensity}), they contain every pure state of $V$
  (\lean{FurtherResults.isSchwartzDensity_proj}), and the von Neumann entropy is finite
  (\lean{FurtherResults.IsSchwartzDensity.vonNeumannEntropy_ne_top}) and continuous on them, by
  Wehrl's theorem (\lean{QuantumMechanics.tendsto_vonNeumannEntropy},
  \lean{FurtherResults.tendsto_vonNeumannEntropy_of_tendsto_numberMoment}). The question about
  the quotient topology remains open.
  (Brief, Sec.~3, QSTO:V; Sec.~5.)
\item \textbf{Field theories.} \emph{Research.} ``For a field theory, we would expect
  the defining observables to be at least countable''; no field-theory example of a quantum
  vector space is given. (Brief, Sec.~5, ``Number of defining observables''.)
\item \textbf{Which observables generate unitary groups.} \emph{Research.} Since
  symmetric and self-adjoint coincide on $V$, ``not all self-adjoint operators generate a unitary
  transformation'', and ``to understand which ones the integrability of the system must be
  checked''. A characterization, analogous to complete vector fields in Hamiltonian mechanics, is
  not given.
  (Brief, Sec.~5, ``Symmetric vs self-adjoint''.)
\item \textbf{Which changes of frame.} \emph{Research.} The proof of Theorem~5 gets [B]WDQ
  from the premise that ``the defining observables of A must be observables in B and
  vice-versa'', which holds exactly when the two frames give the same topology on $V$
  (\lean{topology_eq_iff_continuous}). The brief does not say which changes of frame satisfy it.
  For Example~3 the change of frame $y = \tan(\frac{\pi}{2}\operatorname{erf} x)$, which the
  brief ``Unphysicality of Hilbert spaces for quantum mechanics'' uses against [B]WDQ for
  QSTO:H, does not: it takes a Gaussian out of $\mathcal{S}(\mathbb{R})$
  (\lean{FurtherResults.not_forall_frameUnitary_mem_range}), and the position of the second frame,
  carried back to the first, is not defined at that Gaussian
  (\lean{FurtherResults.toLp_gaussSchwartz_not_mem_domain}; see \lean{further-results.tex}, Sec.~1). So QSTO:V gives [B]WDQ only for a
  class of changes of frame that map $V$ onto itself, as Carcassi, Calder\'on and Aidala already
  require for Schwartz spaces (arXiv:2308.06669v3, Appendix, Sec.~VIII). The class is not
  specified. The restriction cannot simply be dropped: a wave function on $\mathbb{R}$ has finite position
  variance in every smooth coordinate only if it vanishes outside a bounded set; then, unless it
  is zero, its Fourier transform does not, and the free evolution destroys the cutoff at once.
  This is proved in \lean{briefs/HilbertUnphysicality/further-results.tex}, Sec.~2
  (\lean{HilbertUnphysicality.FurtherResults.not_forall_frames_freeEvolution}), and with the
  brief ``Unphysicality of Hilbert spaces for quantum mechanics''
  (\lean{HilbertUnphysicality.quantum_no_cutoff}). Open:
  which group of changes of frame should be part of the data of a quantum vector space, and how
  it relates to the Lie algebra of the defining observables. For plain changes of variable
  $f \mapsto f \circ g$ on $\mathcal{S}(\mathbb{R})$, those that map $\mathcal{S}(\mathbb{R})$ into
  itself are characterized by Galbis and Jord\'a (Rev.\ Mat.\ Iberoam.\ 34 (2018), Thm.~2.3); the
  unitary changes of frame carry the extra factor $\sqrt{g'}$.
  \emph{Partial answer} (\lean{further-results.tex}, Sec.~3). With a group $F$ of changes of
  frame acting on $V$ by linear bijections that preserve the inner product, the premise holds for
  every pair of frames in $F$ if and only if every element of $F$ acts continuously on $V$
  (\lean{FurtherResults.isFrameCompatible_iff_continuous}). Then every change of frame in $F$ is a
  unitary equivalence (\lean{FurtherResults.IsFrameCompatible.unitaryEquiv}). It suffices that
  the transformed defining observables lie in the algebra of the original ones
  (\lean{FurtherResults.isFrameCompatible_of_mem_algebra}); the translations of Example~3 are an
  instance (\lean{FurtherResults.isFrameCompatible_translate}). The choice of $F$ remains open.
  (Brief, Theorem~5, proof; Log, ``Problems found in the brief'', last item.)
\end{enumerate}

\section*{From the formalization}
None. Definition~1 and Examples~3 (in $n$ dimensions) and~4 are checked in full.

\end{document}
